\subsection{Why Daniell closedness is imposed on the codomain}\label{sec:why-complete}

The main theorem (Theorem~\ref{thm:main-theorem}) is stated after restricting the
codomain of the factorization to the category $\mathbf{DInt}$ of Daniell closed
integration systems.
In this subsection we show that this restriction cannot be dispensed with.
That is, the functor $\iota_M\circ\mathcal{D}:\mathbf{Int}\to\mathbf{MInt}$
defined by the Daniell extension is \emph{not} left adjoint to the inclusion
functor $\iota:\mathbf{MInt}\to\mathbf{Int}$.
We then show that $\iota$ does nevertheless have a left adjoint, but that its
reflector is \emph{not} the Daniell extension.

\subsubsection{The Daniell extension is not left adjoint to $\iota$}

\begin{remark}\label{rem:D-I-not-left-adjoint}
  By Lemma~\ref{lem:daniell-ext-is-functor-2}, the Daniell extension defines a
  functor
  \[
    \iota_M\circ\mathcal{D}:\mathbf{Int}\to\mathbf{MInt}
  \]
  from $\mathbf{Int}$ to $\mathbf{MInt}$, but it is not left adjoint to the
  inclusion functor
  \[
    \iota:\mathbf{MInt}\to\mathbf{Int} .
  \]

  Indeed, let $X=\{0,1,2\}$, write $\mathbb{R}^X$ in coordinates, and put
  \[
    L:=\{(a,b,b)\mid a,b\in\mathbb{R}\},
    \qquad
    I(a,b,b):=a .
  \]
  Then $L$ is a Riesz space on $X$ and $I$ is an integral. Moreover, if a
  sequence $\{g_n\}$ of functions in $L$ converges pointwise and monotonically to
  a function $g$, then $g$ too takes equal values at the points $1$ and $2$, and
  $I(g_n)=g_n(0)\to g(0)=I(g)$. Hence $(X,L,I)$ satisfies the monotone
  convergence theorem.

  On the other hand, let $u:=(0,1,1)\in L$. For an arbitrary
  $h=(a,b,c)\in\mathbb{R}^X$, putting
  \[
    f_1=(a,a,a),\quad f_n=0\ (n\ge2),
    \qquad
    \varphi_1=|a|(1,1,1),\quad \varphi_n=u\ (n\ge2) ,
  \]
  we get
  \[
    |f_n|\le\varphi_n,
    \qquad
    \sum_{n=1}^{\infty}I(\varphi_n)=|a|<+\infty.
  \]
  Furthermore $\sum_n\varphi_n(x)<+\infty$ holds only at $x=0$, and at that
  point $h(0)=a=\sum_n f_n(0)$. Hence $h\in L^1$ and
  \[
    L^1=\mathbb{R}^X,
    \qquad
    I^1(h)=h(0).
  \]

  Note here that a morphism of integration systems is determined by its map of
  base spaces. Putting $A:=(X,L,I)$, the condition on a map $\tau:X\to X$ for it
  to define a morphism $A\to A$ is
  \[
    \tau(0)=0,
    \qquad
    \tau(1)=\tau(2)\ \text{or}\ \tau(1),\tau(2)\in\{1,2\} ,
  \]
  so that
  \[
    \left|\operatorname{Hom}_{\mathbf{Int}}(A,\iota A)\right|=5.
  \]
  By contrast, a map $\tau:X\to X$ defining a morphism
  $(\iota_M\circ\mathcal{D})(A)\to A$ must satisfy $h\circ\tau\in L$ for every
  $h\in\mathbb{R}^X$. Hence
  \[
    \tau(0)=0,
    \qquad
    \tau(1)=\tau(2),
  \]
  and
  \[
    \left|\operatorname{Hom}_{\mathbf{MInt}}((\iota_M\circ\mathcal{D})(A),A)\right|=3.
  \]
  Therefore there is no bijection between these two Hom-sets, and
  $\iota_M\circ\mathcal{D}$ is not left adjoint to $\iota$. In this example the
  Daniell extension adjoins all functions on the points $1,2$ that $I$ does not
  distinguish, while the codomain $A$ does not contain them. Hence, in order to
  obtain the universal property of a left adjoint from the Daniell extension
  itself, one must restrict the codomain to the Daniell closed integration
  systems.
\end{remark}

\subsubsection{Another reflection possessed by $\iota$}

That the Daniell extension is not left adjoint to $\iota$ does not mean that
$\iota$ has no left adjoint.
In fact $\iota$ does have a left adjoint.
Its reflector is not the Daniell extension, however, but a different
construction, which takes inside $L^1$ the smallest Riesz subspace satisfying
the monotone convergence theorem.

\begin{theorem}\label{thm:mint-inclusion-has-left-adjoint}
  The inclusion functor
  \[
    \iota:\mathbf{MInt}\to\mathbf{Int}
  \]
  has a left adjoint.
\end{theorem}
\begin{proof}
  Let $A=(X,L,I)$ be an object of $\mathbf{Int}$ and let $(X,L^1,I^1)$ be its
  Daniell extension. Consider the family of subspaces of $L^1$
  \[
    \mathfrak{N}_A
    :=\left\{
      N\subset L^1
      \mathrel{}\middle|\mathrel{}
      \begin{array}{l}
        L\subset N,\quad N\text{ is a Riesz space on $X$},\\
        (X,N,I^1|_N)\in\mathbf{MInt}
      \end{array}
    \right\}
  \]
  and put
  \[
    L^{\mathrm{M}}:=\bigcap_{N\in\mathfrak{N}_A}N,
    \qquad
    I^{\mathrm{M}}:=I^1|_{L^{\mathrm{M}}} .
  \]
  By Lemma~\ref{lem:daniell-ext-contains-source} and
  Lemma~\ref{lem:daniell-ext-is-complete} we have $L\subset L^1$ and
  $(X,L^1,I^1)\in\mathbf{MInt}$, so $L^1\in\mathfrak{N}_A$ and
  $\mathfrak{N}_A$ is nonempty.

  $L^{\mathrm{M}}$ is an intersection of Riesz spaces, hence a Riesz space, and
  $L\subset L^{\mathrm{M}}$. Moreover $I^{\mathrm{M}}$ is a restriction of the
  integral $I^1$, so $(X,L^{\mathrm{M}},I^{\mathrm{M}})$ is an integration
  system.

  We show that this integration system satisfies the monotone convergence
  theorem. Suppose $\{f_n\}\subset L^{\mathrm{M}}$ and $f:X\to\mathbb{R}$ satisfy
  \[
    f_n(x)\uparrow f(x)\quad(\forall x\in X),
    \qquad
    \sup_n I^{\mathrm{M}}(f_n)<+\infty .
  \]
  For every $N\in\mathfrak{N}_A$ we have $f_n\in N$, and $(X,N,I^1|_N)$
  satisfies the monotone convergence theorem, so
  \[
    f\in N,
    \qquad
    I^1(f)=\lim_{n\to\infty}I^1(f_n).
  \]
  As $N$ was arbitrary, $f\in L^{\mathrm{M}}$ and
  \[
    I^{\mathrm{M}}(f)=\lim_{n\to\infty}I^{\mathrm{M}}(f_n).
  \]
  Hence
  \[
    A^{\mathrm{M}}:=(X,L^{\mathrm{M}},I^{\mathrm{M}})
  \]
  is an object of $\mathbf{MInt}$. Since $I^1|_L=I$ by
  Lemma~\ref{lem:daniell-ext-contains-source}, the identity map
  $\operatorname{id}_X$ induces a morphism of integration systems
  \[
    i_A:A\to\iota(A^{\mathrm{M}}) .
  \]

  Let $B=(Y,M,J)$ be an object of $\mathbf{MInt}$ and take a morphism of
  integration systems
  \[
    F:A\to\iota(B) .
  \]
  By Lemma~\ref{lem:daniell-ext-is-functor-1}, the same map of base spaces
  $\underline{F}:Y\to X$ induces a morphism of integration systems
  \[
    F^1:(X,L^1,I^1)\to(Y,M^1,J^1) .
  \]
  Put
  \[
    N_F:=\{f\in L^1\mid F^*f\in M\} .
  \]
  Since $F$ is a morphism of integration systems we have $L\subset N_F$, and
  since $M$ is a Riesz space so is $N_F$. Moreover, for every $f\in N_F$, since
  $F^1$ preserves the integral and $I^1,J^1$ extend the original integrals,
  \[
    J(F^*f)=J^1(F^*f)=I^1(f).
  \]

  We show that $(X,N_F,I^1|_{N_F})$ satisfies the monotone convergence theorem.
  Suppose $\{f_n\}\subset N_F$ and $f:X\to\mathbb{R}$ satisfy
  \[
    f_n(x)\uparrow f(x)\quad(\forall x\in X),
    \qquad
    \sup_n I^1(f_n)<+\infty .
  \]
  Since $(X,L^1,I^1)$ satisfies the monotone convergence theorem,
  \[
    f\in L^1,
    \qquad
    I^1(f)=\lim_{n\to\infty}I^1(f_n).
  \]
  On the other hand $F^*f_n\uparrow F^*f$ and
  \[
    \sup_nJ(F^*f_n)=\sup_nI^1(f_n)<+\infty.
  \]
  Since $B$ satisfies the monotone convergence theorem we get $F^*f\in M$, hence
  $f\in N_F$. Therefore $N_F\in\mathfrak{N}_A$, and we obtain
  $L^{\mathrm{M}}\subset N_F$.

  Hence $\underline{F}:Y\to X$ induces a morphism of integration systems
  \[
    F^\sharp:A^{\mathrm{M}}\to B .
  \]
  Since the map of base spaces underlying $i_A$ is $\operatorname{id}_X$,
  \[
    F=\iota(F^\sharp)\circ i_A.
  \]
  Moreover, if a morphism $H:A^{\mathrm{M}}\to B$ satisfies this equality, then
  comparing the maps of base spaces gives $\underline{H}=\underline{F}$. A
  morphism of integration systems is precisely its map of base spaces, so
  $H=F^\sharp$. Hence $i_A$ is a universal morphism from $A$ to $\iota$.

  Since $A$ was arbitrary,
  Theorem~\ref{thm:left-adjunction-existence-from-universality} shows that
  $\iota$ has a left adjoint.
\end{proof}

\subsubsection{When do the two reflections agree?}

We have thus obtained two reflections for an object $A=(X,L,I)$ of
$\mathbf{Int}$.
\begin{itemize}
  \item the reflection $\mathcal{D}A=(X,L^1,I^1)$ into $\mathbf{DInt}$
    (Theorem~\ref{thm:main-theorem});
  \item the reflection $A^{\mathrm{M}}=(X,L^{\mathrm{M}},I^{\mathrm{M}})$ into
    $\mathbf{MInt}$ (Theorem~\ref{thm:mint-inclusion-has-left-adjoint}).
\end{itemize}
In both cases the base space remains $X$, and by construction
\[
  L\subset L^{\mathrm{M}}\subset L^1 .
\]
When do these two agree?

\begin{lemma}\label{lem:two-reflections-same-daniell}
  Let $A$ be an object of $\mathbf{Int}$. Then
  \[
    \mathcal{D}(A^{\mathrm{M}})=\mathcal{D}(A),
  \]
  that is, $(L^{\mathrm{M}})^1=L^1$, and the extended integrals agree on them.
\end{lemma}
\begin{proof}
  Since $L\subset L^{\mathrm{M}}$ and $I^{\mathrm{M}}|_L=I$, a series
  representation of $f\in L^1$ with respect to $(X,L,I)$ is also, as it stands, a
  series representation with respect to
  $(X,L^{\mathrm{M}},I^{\mathrm{M}})$.
  Hence $L^1\subset(L^{\mathrm{M}})^1$.

  Conversely, since $L^{\mathrm{M}}\subset L^1$ and
  $I^1|_{L^{\mathrm{M}}}=I^{\mathrm{M}}$, the same reasoning gives
  $(L^{\mathrm{M}})^1\subset(L^1)^1$, and
  $(L^1)^1=L^1$ by Lemma~\ref{lem:daniell-ext-is-complete}.
  Together these give $(L^{\mathrm{M}})^1=L^1$.
  That the extended integrals agree is immediate from the definition via series
  representations.
\end{proof}

\begin{proposition}\label{prop:two-reflections-coincide}
  Let $A$ be an object of $\mathbf{Int}$. The following are equivalent.
  \begin{enumerate}
    \item[(i)] $L^{\mathrm{M}}=L^1$, that is, $A^{\mathrm{M}}=\mathcal{D}A$.
    \item[(ii)] $A^{\mathrm{M}}$ is Daniell closed.
    \item[(iii)] The morphism $A^{\mathrm{M}}\to\mathcal{D}A$ induced by the
      identity map $\operatorname{id}_X$ is an isomorphism.
    \item[(iv)] There is no proper Riesz subspace of $L^1$ containing $L$ that
      defines an integration system satisfying the monotone convergence theorem.
  \end{enumerate}
\end{proposition}
\begin{proof}
  (i)$\Rightarrow$(ii) holds because $\mathcal{D}A$ is Daniell closed
  (Lemma~\ref{lem:daniell-ext-is-complete}).
  (ii)$\Rightarrow$(i) holds because, if $A^{\mathrm{M}}$ is Daniell closed, then
  $(L^{\mathrm{M}})^1=L^{\mathrm{M}}$, and the left-hand side equals $L^1$ by
  Lemma~\ref{lem:two-reflections-same-daniell}.
  (i)$\Leftrightarrow$(iii) is immediate, since the objects share the base space
  and a morphism is precisely its map of base spaces.
  (i)$\Leftrightarrow$(iv) holds because $L^{\mathrm{M}}$ is the intersection of
  $\mathfrak{N}_A$, that is, the smallest of the subspaces satisfying the
  conditions.
\end{proof}

\begin{example}
  In the example of Remark~\ref{rem:D-I-not-left-adjoint}, $A=(X,L,I)$ was
  already an object of $\mathbf{MInt}$, so $L^{\mathrm{M}}=L$, while
  $L^1=\mathbb{R}^X$.
  Hence $L^{\mathrm{M}}\subsetneq L^1$ and the two reflections do not agree.

  It is easy to see what is going on in this example.
  The integral $I$ does not distinguish the point $1$ from the point $2$, and the
  elements of $L$ take the same value at these two points.
  The Daniell extension adjoins all functions on these two points, but they are
  not needed in order to satisfy the monotone convergence theorem.
\end{example}

The example above shows that failure of $L$ to separate the points of the base
space can produce a gap between the two reflections.
But separation is not a sufficient condition.
Below we determine $L^{\mathrm{M}}$ in terms of the $\sigma$-algebra generated by
$L$ and derive that fact as a corollary.

\begin{definition}\label{def:sigma-L}
  Let $(X,L,I)$ be an integration system. We call the $\sigma$-algebra on $X$
  \[
    \sigma(L)
    :=\sigma\bigl(\{f^{-1}(B)\mid f\in L,\ B\in\mathcal{B}(\mathbb{R})\}\bigr)
  \]
  the \textbf{$\sigma$-algebra generated by $L$}.
  That is, $\sigma(L)$ is the smallest $\sigma$-algebra making every element of
  $L$ measurable.
  We also put
  \[
    L^{\sigma}:=\{f\in L^1\mid f\text{ is }\sigma(L)\text{-measurable}\},
    \qquad
    I^{\sigma}:=I^1|_{L^{\sigma}} .
  \]
\end{definition}

\begin{proposition}\label{prop:LM-in-Lsigma}
  For every integration system $A=(X,L,I)$,
  \[
    L\subset L^{\mathrm{M}}\subset L^{\sigma}\subset L^1 .
  \]
\end{proposition}
\begin{proof}
  It suffices to show $(X,L^{\sigma},I^{\sigma})\in\mathfrak{N}_A$.

  Every element of $L$ is $\sigma(L)$-measurable (by definition) and
  $L\subset L^1$, so $L\subset L^{\sigma}$.
  The set of all $\sigma(L)$-measurable real-valued functions is a Riesz space on
  $X$, and $L^1$ is a Riesz space as well
  (Theorem~\ref{thm:l1-is-riesz-space}), so their intersection $L^{\sigma}$ is a
  Riesz space on $X$.
  $I^{\sigma}$ is the restriction of the integral $I^1$ to a Riesz subspace, so
  $(X,L^{\sigma},I^{\sigma})$ is an integration system.

  We check the monotone convergence theorem. Suppose
  $\{f_n\}\subset L^{\sigma}$ and $f:X\to\mathbb{R}$ satisfy
  $f_n(x)\uparrow f(x)$ for all $x\in X$ and
  $\sup_nI^{\sigma}(f_n)<+\infty$.
  Since $(X,L^1,I^1)$ satisfies the monotone convergence theorem
  (Lemma~\ref{lem:daniell-ext-is-complete}), we get $f\in L^1$ and
  $I^1(f)=\lim_nI^1(f_n)$.
  Moreover $f$ is a pointwise limit of $\sigma(L)$-measurable functions, hence
  $\sigma(L)$-measurable.
  Therefore $f\in L^{\sigma}$ and $I^{\sigma}(f)=\lim_nI^{\sigma}(f_n)$.

  Hence $L^{\sigma}\in\mathfrak{N}_A$, and since $L^{\mathrm{M}}$ is the
  intersection of $\mathfrak{N}_A$ we obtain
  $L^{\mathrm{M}}\subset L^{\sigma}$.
\end{proof}

\begin{definition}\label{def:supported-on-null}
  Let $(X,L,I)$ be an integration system. We put
  \[
    \mathcal{Z}_I:=\bigl\{h:X\to\mathbb{R}\;\bigm|\;
      \text{there is an $I$-null set $Z$ with $h|_{X\setminus Z}=0$}\bigr\}
  \]
  and call the elements of $\mathcal{Z}_I$ \textbf{functions supported on a null
  set}.
\end{definition}

\noindent
By Lemma~\ref{lem:null-sets-countable-union}, $\mathcal{Z}_I$ is a Riesz space on
$X$.
Indeed, if $h_1$ and $h_2$ vanish outside $Z_1$ and $Z_2$ respectively, then
both $h_1+h_2$ and $|h_1|$ vanish outside $Z_1\cup Z_2$.

\begin{proposition}\label{prop:l1-decomposition}
  For every integration system $A=(X,L,I)$,
  \[
    L^1=L^{\sigma}+\mathcal{Z}_I .
  \]
  Moreover $h\in\mathcal{Z}_I$ implies $I^1(|h|)=0$, so for a decomposition
  $f=g+h$ with $g\in L^{\sigma}$ and $h\in\mathcal{Z}_I$ we have
  $I^1(f)=I^1(g)$.
\end{proposition}
\begin{proof}
  First we check $\mathcal{Z}_I\subset L^1$.
  Let $h\in\mathcal{Z}_I$ and take an $I$-null set $Z$ and a corresponding
  sequence $\{g_n\}$.
  Putting $f_n:=0$ and $\varphi_n:=g_n$ we have $|f_n|\le\varphi_n$ and
  $\sum_nI(\varphi_n)<+\infty$, and any $x$ with $\sum_n\varphi_n(x)<+\infty$
  does not belong to $Z$, so $h(x)=0=\sum_nf_n(x)$.
  Hence $\{f_n\},\{\varphi_n\}$ is a series representation of $h$ and
  $h\in L^1$.
  Moreover $\{g_n\}\subset L\subset L^1$ and $I^1|_L=I$
  (Lemma~\ref{lem:daniell-ext-contains-source}), so $Z$ is a null set for the
  integration system $(X,L^1,I^1)$ as well.
  Since $|h|$ vanishes outside $Z$, applying
  Lemma~\ref{lem:vanish-off-null-set} to $(X,L^1,I^1)$ gives $I^1(|h|)=0$.

  $L^{\sigma}\subset L^1$ holds by definition. Since $L^1$ is a Riesz space
  (Theorem~\ref{thm:l1-is-riesz-space}) we get
  $L^{\sigma}+\mathcal{Z}_I\subset L^1$.

  We show the converse. Let $f\in L^1$, choose a series representation
  $f\overset{\varphi_n}{\simeq}\sum f_n$, and put
  \[
    \Phi_N:=\sum_{n\le N}\varphi_n,
    \qquad
    Z:=\Bigl\{x\in X\;\Bigm|\;\sum_n\varphi_n(x)=+\infty\Bigr\} .
  \]
  Then $\Phi_N\in L$ is $\sigma(L)$-measurable and
  \[
    Z=\bigcap_{M\ge1}\bigcup_{N\ge1}\{\Phi_N>M\} ,
  \]
  so $Z\in\sigma(L)$.
  Moreover $|f_n|\le\varphi_n$ gives $\varphi_n\ge0$, and
  $\sum_nI(\varphi_n)<+\infty$, so $\{\varphi_n\}$ shows that $Z$ is an $I$-null
  set.

  Define a function $g:X\to\mathbb{R}$ by
  \[
    g(x):=
    \begin{cases}
      \sum_nf_n(x) & (x\notin Z)\\
      0 & (x\in Z).
    \end{cases}
  \]
  For $x\notin Z$ the series $\sum_nf_n(x)$ converges absolutely
  (Lemma~\ref{lem:series-exp-is-abs-conv}), so $g$ is real-valued.
  Moreover $g$ is a pointwise limit of the $\sigma(L)$-measurable functions
  $1_{X\setminus Z}\sum_{n\le N}f_n$, hence $\sigma(L)$-measurable.

  The pair $\{f_n\},\{\varphi_n\}$ is also a series representation of $g$.
  Indeed, conditions (1) and (2) are unchanged, and condition (3) is precisely
  the statement that $g(x)=\sum_nf_n(x)$ when
  $\sum_n\varphi_n(x)<+\infty$, that is, when $x\notin Z$.
  Hence $g\in L^1$, and so $g\in L^{\sigma}$.

  Finally, if $x\notin Z$ then $f(x)=\sum_nf_n(x)=g(x)$, so $f-g$ vanishes
  outside $Z$. Therefore $f-g\in\mathcal{Z}_I$ and
  $f=g+(f-g)\in L^{\sigma}+\mathcal{Z}_I$.
\end{proof}

\noindent
The reverse inclusion holds under the Stone condition
(Definition~\ref{def:stone-condition}).
The Stone condition is weaker than requiring that $L$ contain the constant
function $1$, and it includes integration systems of infinite total mass
(Example~\ref{ex:step-functions-R}).

\begin{theorem}\label{thm:LM-equals-Lsigma}
  Let $A=(X,L,I)$ be an integration system satisfying the Stone condition. Then
  \[
    L^{\mathrm{M}}=L^{\sigma}.
  \]
\end{theorem}
\begin{proof}
  Since $L^{\mathrm{M}}\subset L^{\sigma}$ by
  Proposition~\ref{prop:LM-in-Lsigma}, it suffices to show that
  $L^{\sigma}\subset N$ holds for every $N\in\mathfrak{N}_A$.
  Fix $N\in\mathfrak{N}_A$ in what follows.

  \textbf{Step 1: if $g\in L$, $g\ge0$ and $a>0$ then
  $1_{\{g>a\}}\in N\cap L^1$ and $I^1(1_{\{g>a\}})\le I(g)/a$.}
  By the Stone condition $g\wedge a=a\bigl((g/a)\land1\bigr)\in L$, so
  \[
    h_n:=\bigl(n(g-g\land a)\bigr)\land1\in L .
  \]
  Putting $t:=g(x)$ at a point $x$: if $t\le a$ then $t-(t\land a)=0$ and so
  $h_n(x)=0$, and if $t>a$ then $n(t-a)\uparrow+\infty$ and so
  $h_n(x)\uparrow1$.
  Hence $h_n\uparrow1_{\{g>a\}}$ pointwise and monotonically nondecreasing.
  Moreover $0\le h_n\le1$ and $h_n$ vanishes outside $\{g>a\}$, while on
  $\{g>a\}$ we have $g/a>1$, so $h_n\le g/a$ and
  \[
    I(h_n)\le I(g)/a<+\infty.
  \]
  Since $N$ contains $L$ and satisfies the monotone convergence theorem,
  $1_{\{g>a\}}\in N$.
  As $N\subset L^1$ and $I^1$ extends $I$, the estimate on the integral follows
  as well.

  \textbf{Step 2: if $G\in\sigma(L)$ satisfies $1_G\in N$, then $1_E\in N$ for
  every $E\in\sigma(L)$ with $E\subset G$.}
  Put
  \[
    \mathcal{A}_G:=\{E\in\sigma(L)\mid E\subset G,\ 1_E\in N\} .
  \]
  We check that this is a $\sigma$-algebra on $G$.
  We have $G\in\mathcal{A}_G$. If $E\in\mathcal{A}_G$ then
  $1_{G\setminus E}=1_G-1_E\in N$.
  If $E_1,E_2\in\mathcal{A}_G$ then
  $1_{E_1\cup E_2}=1_{E_1}\lor1_{E_2}\in N$.
  Finally, let $E_k\in\mathcal{A}_G$, put $E:=\bigcup_kE_k$ and
  $E'_m:=E_1\cup\dots\cup E_m$; then $E'_m\in\mathcal{A}_G$ and
  $1_{E'_m}\uparrow1_E$, and $E'_m\subset G$ gives
  \[
    I^1(1_{E'_m})\le I^1(1_G)<+\infty .
  \]
  By the monotone convergence theorem for $N$ we get $1_E\in N$.

  Next we show that the trace $\{E\cap G\mid E\in\sigma(L)\}$ of $\sigma(L)$ on
  $G$ is contained in $\mathcal{A}_G$.
  The trace $\sigma$-algebra is generated by the traces of a generating family of
  $\sigma(L)$, so it suffices to check $\{f>a\}\cap G\in\mathcal{A}_G$ for
  $f\in L$ and $a\in\mathbb{R}$.
  \begin{itemize}
    \item Case $a>0$. We have $\{f>a\}=\{f^+>a\}$ with $f^+\in L$ and
      $f^+\ge0$, so Step 1 gives $1_{\{f>a\}}\in N$. Hence
      $1_{\{f>a\}\cap G}=1_{\{f>a\}}\land1_G\in N$.
    \item Case $a=0$. Here $\{f>0\}=\bigcup_{k\ge1}\{f>1/k\}$ is a countable
      union of sets from the previous case.
    \item Case $a<0$. Then $-a>0$, and for $k>1/(-a)$ we have $-a-1/k>0$, so
      \[
        \{f\le a\}=\{-f\ge-a\}=\bigcap_{k>1/(-a)}\{-f>-a-\tfrac1k\}
      \]
      is a countable intersection of sets from the case $a>0$. The claim follows
      from $\{f>a\}\cap G=G\setminus(\{f\le a\}\cap G)$.
  \end{itemize}
  Since $\mathcal{A}_G$ is a $\sigma$-algebra it contains the whole trace.
  An $E\in\sigma(L)$ with $E\subset G$ belongs to the trace as $E=E\cap G$, so
  $1_E\in N$.

  \textbf{Step 3: $L^{\sigma}\subset N$.}
  Both $L^{\sigma}$ and $N$ are Riesz spaces and $f$ decomposes as
  $f=f^+-f^-$, so it suffices to treat the case $f\in L^{\sigma}$ with $f\ge0$.
  Choose a series representation $f\overset{\varphi_n}{\simeq}\sum f_n$.

  At a point $x$ where $\varphi_n(x)=0$ for all $n$ we have
  $\sum_n\varphi_n(x)=0<+\infty$, so $f(x)=\sum_nf_n(x)$, and
  $|f_n(x)|\le\varphi_n(x)=0$ gives $f(x)=0$. Hence
  \[
    \{f\ne0\}\subset\bigcup_{n,k\ge1}G_{n,k},
    \qquad
    G_{n,k}:=\{\varphi_n>1/k\} .
  \]
  Since $\varphi_n\in L$ and $\varphi_n\ge0$, Step 1 gives
  $1_{G_{n,k}}\in N\cap L^1$, and $G_{n,k}\in\sigma(L)$.
  Arranging the family $\{G_{n,k}\}_{n,k}$ into a single sequence
  $G_1,G_2,\dots$ and putting $H_m:=G_1\cup\dots\cup G_m$, we get
  $1_{H_m}=1_{G_1}\lor\dots\lor1_{G_m}\in N$ and $H_m\in\sigma(L)$.

  Take $c>0$. Since $f$ is $\sigma(L)$-measurable we have
  $\{f\ge c\}\in\sigma(L)$, and $\{f\ge c\}\cap H_m$ belongs to $\sigma(L)$ and
  is contained in $H_m$, so Step 2 gives
  $1_{\{f\ge c\}\cap H_m}\in N$.
  From $c>0$ and $f\ge0$ we get
  $\{f\ge c\}\subset\{f\ne0\}\subset\bigcup_mH_m$, so
  $\{f\ge c\}\cap H_m\uparrow\{f\ge c\}$, that is,
  $1_{\{f\ge c\}\cap H_m}\uparrow1_{\{f\ge c\}}$.
  Moreover $f/c\ge1$ on $\{f\ge c\}$ and $f/c\ge0$ off it, so
  $1_{\{f\ge c\}}\le f/c$ and
  \[
    I^1(1_{\{f\ge c\}\cap H_m})\le I^1(f)/c<+\infty.
  \]
  By the monotone convergence theorem for $N$ we get $1_{\{f\ge c\}}\in N$.

  Finally, putting
  \[
    f_m:=\sum_{k=1}^{m2^m}2^{-m}1_{\{f\ge k2^{-m}\}}
    \qquad\Bigl(=\bigl(2^{-m}\lfloor2^mf\rfloor\bigr)\land m\Bigr) ,
  \]
  the function $f_m$ belongs to $N$ as a finite linear combination of elements of
  $N$.
  It is nondecreasing in $m$ and converges pointwise to $f$, and
  $0\le f_m\le f$ gives $I^1(f_m)\le I^1(f)<+\infty$.
  By the monotone convergence theorem for $N$ we get $f\in N$.

  Since $N\in\mathfrak{N}_A$ was arbitrary we have
  $L^{\sigma}\subset L^{\mathrm{M}}$, and together with
  Proposition~\ref{prop:LM-in-Lsigma} we obtain
  $L^{\mathrm{M}}=L^{\sigma}$.
\end{proof}

\begin{corollary}\label{cor:LM-equals-Lsigma-one}
  If the constant function $1$ belongs to $L$ then
  $L^{\mathrm{M}}=L^{\sigma}$.
\end{corollary}
\begin{proof}
  Since $1\in L$ and $L$ is a Riesz space, $f\land1\in L$ for every $f\in L$,
  that is, the Stone condition holds.
\end{proof}

\noindent
When $1\in L$ one may take $G:=X$, so the localization to $H_m$ in Step 2 of the
proof above is unnecessary.

\medskip
\noindent
\textbf{Relation to classical Daniell--Stone theory.}
The content of Theorem~\ref{thm:LM-equals-Lsigma} lies in the neighbourhood of
the classical theory.
The smallest family containing a Riesz space satisfying the Stone condition (a
Stone vector lattice) and closed under monotone limits is a standard object,
called the $\sigma$-ring of real functions generated by $L$, and that it
corresponds to $\sigma(L)$-measurability is classically known.
\S4 of Chapter 16 of \cite{Royden}, ``Measurability and measure'', is this part.

What distinguishes the present theorem from the classical assertion is that the
monotone limits are taken inside $L^1$, and only under the boundedness condition
$\sup_nI(f_n)<+\infty$ on the integrals
(Definition~\ref{def:satisfies-mct}).
This is why the generated monotone family cannot be used as it stands, and it is
also the source of the truncation in Step 1 and the localization to $H_m$ in the
proof.
The conclusion obtained is the identification of \textbf{that restricted closure
with the second reflection $L^{\mathrm{M}}$}
(Theorem~\ref{thm:mint-inclusion-has-left-adjoint}).

\begin{corollary}\label{cor:l1-decomposition-LM}
  If $A=(X,L,I)$ satisfies the Stone condition then
  \[
    L^1=L^{\mathrm{M}}+\mathcal{Z}_I .
  \]
  That is, \textbf{the difference between $L^1$ and the second reflection
  $L^{\mathrm{M}}$ amounts exactly to the functions supported on null sets}.
\end{corollary}
\begin{proof}
  By Theorem~\ref{thm:LM-equals-Lsigma} we have
  $L^{\mathrm{M}}=L^{\sigma}$, so the claim follows from
  Proposition~\ref{prop:l1-decomposition}.
\end{proof}

\begin{remark}[two-stage factorization of the Daniell extension]\label{rem:D-factors}
  Under the Stone condition, the results so far split the Daniell extension into
  the following two stages.
  \[
    L
    \quad\subset\quad
    \underbrace{L^{\mathrm{M}}=L^{\sigma}}_{\text{first stage}}
    \quad\subset\quad
    \underbrace{L^1=L^{\mathrm{M}}+\mathcal{Z}_I}_{\text{second stage}}
  \]
  The first stage is the operation of closing $L$ under monotone limits, and it
  is obtained as the reflection into $\mathbf{MInt}$
  (Theorem~\ref{thm:mint-inclusion-has-left-adjoint}).
  Its destination is the set of all $\sigma(L)$-measurable integrable functions
  (Theorem~\ref{thm:LM-equals-Lsigma}).
  The second stage is the operation of adjoining the functions supported on null
  sets (Corollary~\ref{cor:l1-decomposition-LM}).

  Read in the language of measures, the first stage is the operation of taking
  the integrable functions measurable with respect to the $\sigma$-algebra
  generated by $L$, and the second is the operation of completing the measure
  (passing through Theorem~\ref{thm:daniell-stone}).
  That is,
  \begin{quote}
    \textbf{the Daniell extension is the composite of Daniell--Stone style
    measurability and completion of the measure.}
  \end{quote}

  \noindent
  Both stages are classical operations.
  What this paper supplies is the identification of this factorization
  (Theorem~\ref{thm:LM-equals-Lsigma} and
  Proposition~\ref{prop:l1-decomposition}) and the categorical placement of the
  composite as a reflector.
  As stated in Subsection~\ref{sec:prior-work},
  \textbf{the position of this paper, in a phrase, is ``the Daniell counterpart
  of the completion of a measure space''.}
\end{remark}

\begin{corollary}\label{cor:when-coincide-sigma}
  Let $A=(X,L,I)$ be an integration system satisfying the Stone condition.
  Then the equivalent conditions of
  Proposition~\ref{prop:two-reflections-coincide} are equivalent to the
  following.
  \begin{enumerate}
    \item[(v)] Every element of $L^1$ is $\sigma(L)$-measurable.
  \end{enumerate}
\end{corollary}
\begin{proof}
  By Theorem~\ref{thm:LM-equals-Lsigma} we have
  $L^{\mathrm{M}}=L^{\sigma}$, and $L^{\sigma}=L^1$ is a restatement of (v).
\end{proof}

\noindent
This is a characterization in terms of properties of $(X,L,I)$ itself, and makes
no reference to the construction of $L^{\mathrm{M}}$.

\begin{corollary}\label{cor:lebesgue-reflections-differ}
  Let $A$ be the integration system of Example~\ref{ex:step-functions}, that is,
  the integration system consisting of the step functions on $[0,1)$ and the
  Riemann integral. Then
  \[
    L^{\mathrm{M}}=\{f:[0,1)\to\mathbb{R}\text{ Borel measurable and Lebesgue integrable}\}
    \subsetneq L^1
    =\{f:[0,1)\to\mathbb{R}\text{ Lebesgue integrable}\} ,
  \]
  and the two reflections do not agree.
\end{corollary}
\begin{proof}
  Truncating a step function at $1$ is again a step function, so this
  integration system satisfies the Stone condition
  (Example~\ref{ex:stone-holds}).
  Moreover $\sigma(L)$ is the $\sigma$-algebra generated by the indicator
  functions of the half-open intervals $[a,b)\subset[0,1)$, that is, the Borel
  $\sigma$-algebra of $[0,1)$.
  The equalities follow from the fact that $L^1$ is the set of all Lebesgue
  integrable functions (Example~\ref{ex:lebesgue}) together with
  Theorem~\ref{thm:LM-equals-Lsigma}.

  We check that the inclusion is proper. The Cantor set $C\subset[0,1]$ is a
  Lebesgue null set and has the cardinality of the continuum $\mathfrak{c}$.
  Hence $C\setminus\{1\}$ has $2^{\mathfrak{c}}$ subsets, while $[0,1)$ has only
  $\mathfrak{c}$ Borel sets.
  There is therefore a non-Borel set $E\subset C\setminus\{1\}$.
  As a subset of a Lebesgue null set, $E$ is Lebesgue measurable, and $1_E$ is
  Lebesgue integrable with integral $0$, so $1_E\in L^1$.
  On the other hand $E\notin\mathcal{B}([0,1))$, so $1_E$ is not Borel
  measurable and $1_E\notin L^{\mathrm{M}}$.
\end{proof}

\begin{remark}\label{rem:separation-not-sufficient}
  For the integration system of
  Corollary~\ref{cor:lebesgue-reflections-differ}, $L$ separates the points of
  $[0,1)$ (for $x\ne y$ one can take a half-open interval containing $x$ and not
  $y$).
  It satisfies the Stone condition as well
  (Example~\ref{ex:stone-holds}).
  Nevertheless the two reflections do not agree.
  Hence \textbf{neither separation of points nor the Stone condition is a
  sufficient condition for the two reflections to agree}.

  What produces the gap can be identified in general, not only in this example.
  By Proposition~\ref{prop:l1-decomposition} we have
  $L^1=L^{\sigma}+\mathcal{Z}_I$ for every integration system, and under the
  Stone condition $L^1=L^{\mathrm{M}}+\mathcal{Z}_I$ by
  Corollary~\ref{cor:l1-decomposition-LM}.
  That is, \textbf{the gap between the two reflections amounts exactly to the
  functions supported on null sets}.
  In this example it appears in the form that $L^{\mathrm{M}}$ contains only
  Borel measurable functions whereas $L^1$ contains functions behaving
  arbitrarily on null sets.
  That is, it is the difference between Borel and Lebesgue, the amount by which
  the measure is completed.

  It is the fact that the Daniell extension imposes condition (3) of a series
  representation ``only at the points where the sum of the dominating functions
  converges'' (see also Remark~\ref{rem:why-half-open}) that performs this
  completion automatically.
\end{remark}

\begin{example}
  Let us verify Theorem~\ref{thm:LM-equals-Lsigma} in the example of
  Remark~\ref{rem:D-I-not-left-adjoint} ($X=\{0,1,2\}$, $L=\{(a,b,b)\}$,
  $I(a,b,b)=a$).
  Here $1=(1,1,1)\in L$.
  The elements of $L$ take the same value at the points $1$ and $2$, so
  \[
    \sigma(L)=\{\emptyset,\{0\},\{1,2\},X\} ,
  \]
  and the $\sigma(L)$-measurable functions are precisely the functions taking the
  same value at the points $1$ and $2$, that is, the elements of $L$.
  Hence $L^{\sigma}=L^1\cap L=L$, and $L^{\mathrm{M}}=L$ as the theorem asserts.
  This agrees with the value computed directly in the example above.
\end{example}

\begin{example}\label{ex:step-functions-R}
  Theorem~\ref{thm:LM-equals-Lsigma} applies to integration systems of infinite
  total mass as well.
  Let $X=\mathbb{R}$, let $L$ be the set of all step functions with bounded
  support, and let $I$ be their integral.
  Then $1_{\mathbb{R}}\notin L$, but truncating a step function at $1$ is again
  a step function, so the Stone condition is satisfied.
  Moreover $\sigma(L)$ is the $\sigma$-algebra generated by the bounded
  half-open intervals, that is, the Borel $\sigma$-algebra of $\mathbb{R}$.

  Hence, by Theorem~\ref{thm:LM-equals-Lsigma}, $L^{\mathrm{M}}$ is the set of
  all Borel measurable elements of $L^1$.
  The two reflections do not agree for this integration system either.
  Indeed, the Cantor set $C\subset[0,1]$ is an $I$-null set (take step functions
  $g_n$ covering $C$ by finitely many intervals of total length at most
  $2^{-n}$), and by the cardinality argument (the same as in the proof of
  Corollary~\ref{cor:lebesgue-reflections-differ}) there is a non-Borel set
  $E\subset C$.
  Since $1_E$ vanishes outside $C$ we have
  $1_E\in\mathcal{Z}_I\subset L^1$
  (Proposition~\ref{prop:l1-decomposition}), but $1_E$ is not Borel measurable,
  so it does not belong to $L^{\mathrm{M}}$.
\end{example}

\begin{problem}\label{prob:when-reflections-coincide}
  Is the hypothesis ``the Stone condition'' in
  Theorem~\ref{thm:LM-equals-Lsigma} necessary?
  That is, is there an integration system that does not satisfy the Stone
  condition and has $L^{\mathrm{M}}\ne L^{\sigma}$?
\end{problem}

\noindent
That the Stone condition is not necessary is shown by
Example~\ref{ex:stone-fails}.
That example ($X=\{1,2\}$, $L=\{(a,2a)\}$, $I(a,2a)=a$) does not satisfy the
Stone condition, but $L^1=L$ as seen there.
Moreover $\sigma(L)=2^X$ from the preimages under $f=(1,2)$, so
$L^{\sigma}=L^1=L$ and therefore $L^{\mathrm{M}}=L^{\sigma}$ holds.
The question is whether the equality can fail once the Stone condition is
dropped.

The Stone condition was used only in Step 1 of the proof, that is, at the point
where the indicator functions of the upper sets of elements of $L$ are extracted
inside $N$.
Without that operation there is no guarantee that the indicator functions of
elements of $\sigma(L)$ belong to $N$, and the footing of Steps 2 and 3 is lost.
